\section{两个实例：Fold6/HPA 与 Zeckendorf 截断折叠}

本文框架是抽象的；但其可审计性与可复现性依赖于把 \(\Clo\) 具体化为可计算对象。本节给出两个实例化入口，并明确它们在本文中的角色是“本体约束的具体选择”。

\subsection{实例 1：\(m=6\) 的无损分块折叠核（Fold6/HPA）}

该实例提供了一个最小可计算的“ontic 可逆 vs epistemic 搜索”模型：当 uplift 边标签可读时，块级逆映射确定；当 uplift 不可见时，逆向成为一致性搜索，其算法成本可直接作为 \(\kappaobs\) 的具体实现。进一步地，由于 uplift 字母表有限且每步记录代价为 \(O(1)\) 比特，本实例也为 \(\kappaobs\sim\tautime\) 的同阶等价提供了具体支撑：时间轴可由 uplift 胶带的编码长度 \(\tautime\) 给出。

\subsection{实例 2：一般 \(m\) 的 Zeckendorf 截断折叠与时间纤维}

该实例强调“可见稳定态 \(F_{m+2}\) 与被截断尾部（时间纤维）”的组织方式：本体层只固定折叠映射与闭包图一致性；观察者层把“回到过去”转化为在纤维结构上的 lifting 搜索与分支约束。若把被截断尾部视为必须补齐的隐藏胶带，则其规模自然给出负信息时间 \(\tautime\) 的一个实例化解释。

\begin{remark}[完备系统实例化（折叠核 + 明确 uplift 规则）]
以 Zeckendorf 截断折叠为核时，可以把“时间/uplift”具体化为窗口外 Zeckendorf 尾部，并用尾部移位（\(\mathrm{Tail}\)）给出明确的 uplift 规则；此时“向上搜索”可被实现为 \(\mathrm{Tail}^{-1}\) 的分支扩展，而能量预算可用并行分支上限（beam width）建模。该方向的系统化整理与可复现实验见 \path{docs/papers/math/2026_zeckendorf_complete_system_energy_uplift/}。
\end{remark}

\subsection{本文实验的实例选择与解释边界}

为了把“Hilbert 尺子改变塌缩速度”做成可复现实验，本文在数值部分采用一种标准化的观察者推断算子：观察者先构造一个由“折叠边 + 选定几何邻接边”组成的辅助图对象，然后在该图上运行 WL1 细化并记录解析时间（resolve\_t）。这里需要强调两点：
\begin{itemize}
  \item 本体层固定折叠映射（以及其诱导的宏微关系）；几何邻接边被视为观察者的先验与推断结构选择，而非被当作物理本体的必然输入。
  \item 因此，本文关于“维度/尺子”的结论应被理解为：在给定折叠核下，不同局部性组织会系统性改变推断过程生成的搜索/细化时空图与成本分布。
\end{itemize}

\begin{remark}
为避免 PDF 书签中的数学标记触发警告，本文正文不在标题中使用 \(\cdot\) 形式的内联数学；必要的数学符号在段落内给出即可。
\end{remark}

